\begin{pastexam}[label={px:ncu-113-dsa-17}]{中央資工所 113 學年度 第 17 題}
（複選；每題 5 分，全部答對才給分，答錯倒扣一分）

Consider a hash table of size 7, with hash function $H(k)=k\ \%\ 7$, and pseudo random $i=(i+5)\ \%\ 7$. We
want to insert the following keys one by one from left to right.

15, 11, 25, 16, 9, 8, 12

If random probing is used, which of the following statements are incorrect?
\begin{enumerate}[label=(\Alph*)]
\item key 25 is at position 4.
\item key 16 is at position 0.
\item key 9 is at position 0.
\item key 8 is at position 6.
\item Some key has no place to store in the table.
\end{enumerate}
\end{pastexam}

\begin{pastsoln}
題目問的是\textbf{錯誤}的敘述。 答案：(A)(C)(E)。

\textbf{探測方式}。 先試 $H(k)$；有人佔了，就用 $i \gets (i + 5) \bmod 7$ 算下一格，也就是每次往後跳
$5$ 格（$i$ 當成位置）。 如果把 $i$ 當成加在 $H(k)$ 上的位移，位移依序是 $0, 5, 10 \bmod 7 = 3, 1, \ldots$，
試的格子是 $H(k) + 5j \bmod 7$，和前一種讀法完全一樣。 $5$ 和 $7$ 互質，所以每個鍵的探測序列會走遍全部
$7$ 格（資結書第~12.3~節），表沒滿就一定插得進去。

\textbf{逐一插入}：
\begin{center}
\small
\begin{tabular}{rcl}
\toprule
鍵 & $k \bmod 7$ & 試過的格子（最後一個是放進去的位置） \\
\midrule
15 & 1 & 1 \\
11 & 4 & 4 \\
25 & 4 & 4（有 11），$(4+5) \bmod 7 = 2$ \\
16 & 2 & 2（有 25），$0$ \\
9 & 2 & 2（有 25），0（有 16），$5$ \\
8 & 1 & 1（有 15），$6$ \\
12 & 5 & 5（有 9），$(5+5) \bmod 7 = 3$ \\
\bottomrule
\end{tabular}
\end{center}
最後的表：
\begin{center}
\small
\begin{tabular}{c*{7}{c}}
\toprule
位置 & 0 & 1 & 2 & 3 & 4 & 5 & 6 \\
\midrule
鍵 & 16 & 15 & 25 & 12 & 11 & 9 & 8 \\
\bottomrule
\end{tabular}
\end{center}

逐項檢查：
\begin{itemize}
\item (A) 錯誤，要選。 $25$ 在第 $2$ 格；第 $4$ 格是 $11$。
\item (B) 正確，不選。
\item (C) 錯誤，要選。 $9$ 在第 $5$ 格；第 $0$ 格是 $16$。
\item (D) 正確，不選。
\item (E) 錯誤，要選。 七個鍵剛好填滿七格，每個鍵都放進去了。
\end{itemize}
\end{pastsoln}
